% Part: first-order-logic
% Chapter: axiomatic-deduction
% Section: proving-things

\documentclass[../../../include/open-logic-section]{subfiles}

\begin{document}

\iftag{FOL}
      {\olfileid{fol}{axd}{pro}}
      {\olfileid{pl}{axd}{pro}}

\olsection{Voorbeelden van \usetoken{P}{derivation}}


\begin{ex}
Stel dat we $(\lnot !D \lor !E) \lif (!D \lif !E)$ willen bewijzen.
Dit is duidelijk geen instantie van een van onze axioma's, dus moeten
we de regel \MP{} gebruiken om de formule !!{derive}. Onze enige regel
is~MP: uit $!A$ en $!A \lif !B$ mogen we~$!B$ als gerechtvaardigde stap
opnemen. Een mogelijke strategie is \olref[prp]{ax:lor3} te gebruiken
met $!A$ gelijk aan $\lnot !D$, $!B$ gelijk aan $!E$ en $!C$ gelijk
aan $!D \lif !E$. Dat geeft de instantie
\[
(\lnot !D \lif (!D \lif !E)) \lif ((!E \lif (!D \lif !E)) \lif ((\lnot
!D \lor !E) \lif (!D \lif !E))).
\]
Waarom? Twee toepassingen van MP leveren het laatste deel op, en dat
is precies wat we willen. Verder zien we direct dat $\lnot !D \lif
(!D \lif !E)$ een instantie is van \olref[prp]{ax:lnot2} en dat $!E
\lif (!D \lif !E)$ een instantie is van \olref[prp]{ax:lif1}. Onze
!!{derivation} luidt dus:
\begin{derivation}
  1. &  $\lnot !D \lif (!D \lif !E)$ & \olref[prp]{ax:lnot2} \\
  2. & $(\lnot !D \lif (!D \lif !E)) \lif {}$\\
  &\qquad $((!E \lif (!D \lif !E)) \lif ((\lnot !D \lor !E) \lif (!D \lif !E)))$ & \olref[prp]{ax:lor3}\\
  3. & $(!E \lif (!D \lif !E)) \lif ((\lnot !D \lor !E) \lif (!D \lif !E))$ & 1, 2, \MP\\
  4. & $!E \lif (!D \lif !E)$ & \olref[prp]{ax:lif1}\\
  5. & $(\lnot !D \lor !E) \lif (!D \lif !E)$ & 3, 4, \MP
\end{derivation}
\end{ex}

\begin{ex}\ollabel{ex:identity}
Laten we proberen !!a{derivation} van $!D \lif !D$ te vinden. Deze
formule is geen instantie van een axioma, dus moeten we \MP{} gebruiken
om haar !!{derive}. \olref[prp]{ax:lif1} is een axioma van de vorm~$!A
\lif !B$ waarop we~\MP kunnen toepassen. Om ons te helpen, moet de
formule~$!B$ die \MP{} hier als correcte stap rechtvaardigt natuurlijk
$!D \lif !D$ zijn, want die formule willen we !!{derive}. Dan moet~$!A$
ook $!D$ zijn. We zouden dus de volgende instantie van
\olref[prp]{ax:lif1} kunnen nemen:
\[
!D \lif (!D \lif !D)
\]
Om \MP toe te passen, moeten we ook de bijbehorende tweede premisse,
namelijk~$!A$, rechtvaardigen. In ons geval is dat echter~$!D$, en
$!D$ op zichzelf zullen we niet kunnen !!{derive}. We hebben dus een
andere strategie nodig.

Het andere axioma waarin alleen~$\lif$ voorkomt, is
\olref[prp]{ax:lif2}, namelijk
\[
(!A \lif (!B \lif !C)) \lif ((!A \lif !B) \lif (!A \lif !C))
\]
Met twee toepassingen van \MP{} kunnen we de laatste geneste
implicatie bereiken. We zoeken dus opnieuw een instantie, ditmaal van
\olref[prp]{ax:lif2}, waarin $!A \lif !C$ gelijk is aan $!D \lif !D$,
de !!{formula} waarop we mikken. Dan zijn zowel~$!A$ als~$!C$ gelijk
aan~$!D$. Hoe kiezen we~$!B$ zo dat zowel $!A \lif (!B \lif !C)$ als
$!A \lif !B$, in ons geval dus $!D \lif (!B \lif !D)$ en $!D \lif
!B$, ook !!{derivable} zijn? De eerste formule is al een instantie van
\olref[prp]{ax:lif1}, ongeacht onze keuze voor~$!B$. En $!D \lif !B$
is eveneens een instantie van \olref[prp]{ax:lif1} als $!B$ gelijk is
aan $(!D \lif !D)$. Onze !!{derivation} is dus:
\begin{derivation}
  1. & $!D \lif ((!D \lif !D) \lif !D)$ & \olref[prp]{ax:lif1}\\
  2. & $(!D \lif ((!D \lif !D) \lif !D)) \lif {}$\\
  & \qquad $((!D \lif (!D \lif !D)) \lif (!D \lif !D))$ & \olref[prp]{ax:lif2}\\
  3. & $(!D \lif (!D \lif !D)) \lif (!D \lif !D)$ & 1, 2, \MP\\
  4. & $!D \lif (!D \lif !D)$ & \olref[prp]{ax:lif1}\\
  5. & $!D \lif !D$ & 3, 4, \MP
\end{derivation}
\end{ex}

\begin{ex}\ollabel{ex:chain}
Soms willen we aantonen dat er !!a{derivation} bestaat van een
!!{formula} uit andere !!{formula}s~$\Gamma$. Laten we bijvoorbeeld
aantonen dat we $!A \lif !C$ kunnen !!{derive} uit $\Gamma = \{!A \lif
!B, !B \lif !C\}$.
\begin{derivation}
  1. & $!A \lif !B$ & \Hyp\\
  2. & $!B \lif !C$ & \Hyp\\
  3. & $(!B \lif !C) \lif (!A \lif (!B \lif !C))$ & \olref[prp]{ax:lif1} \\
  4. & $!A \lif (!B \lif !C)$ & 2, 3, \MP\\
  5. & $(!A \lif (!B \lif !C)) \lif {}$\\
  & \qquad $((!A \lif !B) \lif (!A \lif !C))$ & \olref[prp]{ax:lif2}\\
  6. & $((!A \lif !B) \lif (!A \lif !C))$ & 4, 5, \MP\\
  7. & $!A \lif !C$ & 1, 6, \MP
\end{derivation}
De regels met het label ``\Hyp'' (voor ``hypothese'') geven aan dat de
!!{formula} op die regel !!a{element} is van~$\Gamma$.
\end{ex}

\begin{prop}
  \ollabel{prop:chain} Als $\Gamma \Proves !A \lif !B$ en $\Gamma
  \Proves !B \lif !C$, dan $\Gamma \Proves !A \lif !C$.
\end{prop}

\begin{proof}
  Stel dat $\Gamma \Proves !A \lif !B$ en $\Gamma \Proves !B \lif
  !C$. Dan bestaat er !!a{derivation} van $!A \lif !B$ uit~$\Gamma$,
  en eveneens !!a{derivation} van~$!B \lif !C$ uit~$\Gamma$. Voeg
  deze samen tot één !!{derivation} door ze achter elkaar te plaatsen.
  Voeg vervolgens de regels 3--7 van de !!{derivation} uit het vorige
  voorbeeld toe. Dit is !!a{derivation} van $!A \lif !C$---de laatste
  regel van de nieuwe !!{derivation}---uit~$\Gamma$. De
  rechtvaardigingen van de regels 4 en~7 blijven geldig als we de
  verwijzing naar regel~1 vervangen door een verwijzing naar de
  laatste regel van de !!{derivation} van~$!A \lif !B$, en de
  verwijzing naar regel~2 door een verwijzing naar de laatste regel
  van de !!{derivation} van~$!B \lif !C$.
\end{proof}

\begin{prob} 
Toon de volgende beweringen aan door !!{derivation}s uit de axioma's
te geven:
\begin{enumerate} 
\item $(!A \land !B) \lif (!B \land !A)$
\item $((!A \land !B) \lif !C) \lif (!A \lif (!B \lif !C))$
\item $\lnot(!A \lor !B) \lif \lnot !A$
\end{enumerate} 
\end{prob}



\end{document}
